\section{Light --- Compiler --- Call Frames}
%%% generated by the blueprint pipeline from TCSlib.Programs.Light.Compiler.CallFrames
\subsection*{Light --- Compiler --- Call Frames: Introduction}

Declarations of the Lean module \texttt{TCSlib.Programs.Light.Compiler.CallFrames}.

\subsection*{Light --- Compiler --- Call Frames: Results}

\begin{theorem}[Straight instructions in argument code]
\lean{Light.Compiler.straight_argsCode}
\label{Light.Compiler.straight_argsCode}
\leanok
\uses{Light.Compiler.argsCode, Light.Compiler.straight_compileExpr, Light.Expr, WordRAM.Straight}
\difficulty{4}
For every finite list of expressions and every natural number $k$, each instruction in
the code that compiles the list of expressions starting at $k$ is straight.
\end{theorem}

\begin{theorem}[Compiled argument evaluation]
\lean{Light.Compiler.effects_argsCode}
\label{Light.Compiler.effects_argsCode}
\leanok
\uses{Light.Compiler.Framed, Light.Compiler.Framed.same, Light.Compiler.Temps, Light.Compiler.argsCode, Light.Compiler.cT, Light.Compiler.effects_compileExpr_of_width, Light.Compiler.temps_subset_scratch, Light.Expr, Light.Expr.Safe, Light.Expr.val, Light.Expr.width, Light.frame, WordRAM.AgreeOutside, WordRAM.AgreeOutside.mono, WordRAM.AgreeOutside.read, WordRAM.AgreeOutside.refl, WordRAM.AgreeOutside.trans, WordRAM.effects, WordRAM.effects_append, WordRAM.wd}
\difficulty{5}
Let $Z$ be a word-RAM layout, $\sigma$ a light state, $q$ a stack position, and $m$ a
memory map, and suppose the frame relation $\mathrm{Framed}(Z,\sigma,q,m)$ holds. Let
$es$ be a list of expressions and $k$ a starting temporary index. Suppose each expression
in $es$ has width at most $Z.F$ and is safe in state $\sigma$ under the limits $Z.\mathrm{lim}$,
and $k+|es|\le Z.F+1$. Executing the compiled argument code places the $W$-bit encoding
of the value of the $i$th expression in temporary cell $k+i$ for every $i<|es|$, and
changes memory only in the temporary cells indexed from $k$ through $Z.F$.
\end{theorem}

\begin{theorem}[Successor frame initialization]
\lean{Light.Compiler.frameInit_succ}
\label{Light.Compiler.frameInit_succ}
\leanok
\uses{Light.Compiler.cADDR, Light.Compiler.cNFP, Light.Compiler.cPool, Light.Compiler.cT, Light.Compiler.cZERO, Light.Compiler.frameInit}
\difficulty{1}
For natural numbers $F'$ and $n$, initializing $F'+1$ frames is the concatenation of
initializing the first $F'$ frames with the two instructions that compute the address
using $c_{\mathrm{ADDR}}$ and $c_{\mathrm{NFP}}$ and the pool cell at index $2F'$, then
store either $c_T(F')$ when $F'<n$ or $c_{\mathrm{ZERO}}$ otherwise.
\end{theorem}

\begin{theorem}[Straight instructions in frame initialization]
\lean{Light.Compiler.straight_frameInit}
\label{Light.Compiler.straight_frameInit}
\leanok
\uses{Light.Compiler.frameInit, Light.Compiler.frameInit_succ, WordRAM.Straight}
\difficulty{4}
For every natural number $n$ and every natural number $F'$, each instruction in the
frame-initialization code of size $F'$ for $n$ is straight.
\end{theorem}

\begin{theorem}[Effects of frame initialization]
\lean{Light.Compiler.effects_frameInit}
\label{Light.Compiler.effects_frameInit}
\leanok
\uses{Light.Compiler.Sizes, Light.Compiler.Sizes.inRange_stack, Light.Compiler.Stack, Light.Compiler.cADDR, Light.Compiler.cNFP, Light.Compiler.cPool, Light.Compiler.cStack, Light.Compiler.cStack_sub, Light.Compiler.cT, Light.Compiler.cZERO, Light.Compiler.frameInit, Light.Compiler.frameInit_succ, WordRAM.AgreeOutside, WordRAM.AgreeOutside.cell, WordRAM.AgreeOutside.read, WordRAM.AgreeOutside.refl, WordRAM.AgreeOutside.update, WordRAM.effect_store, WordRAM.effect_sub, WordRAM.effects, WordRAM.effects_append, WordRAM.effects_cons, WordRAM.effects_nil, WordRAM.wd}
\difficulty{6}
Let $Z$ be compiler size data with word size $W$, stack capacity $Q$, and frame size
$F$, and let $n,q',F'$ be natural numbers and $m$ an initial memory. Suppose the
new-frame pointer contains the stack address for $q'$, each pool cell indexed by $2i$
contains the integer $2i$ for every $i<F'$, and $q'+F'\le Q$. After executing the
frame-initialization instructions, each stack cell indexed by $q'+i$, for $i<F'$,
contains the initial value of temporary cell $i$ if $i<n$, and otherwise the initial
value of the zero cell; moreover, memory agrees with its initial value outside the
address scratch cell and the stack cells indexed from $q'$ through $Q-1$.
\end{theorem}

\begin{theorem}[Straight instructions in a new frame]
\lean{Light.Compiler.straight_newFrame}
\label{Light.Compiler.straight_newFrame}
\leanok
\uses{Light.Compiler.newFrame, Light.Compiler.straight_frameInit, WordRAM.Straight}
\difficulty{3}
For all natural numbers $F,n,$ and $r$, every instruction in the code that creates a new
frame of size $F$ for $n$ arguments with return index $r$ is straight.
\end{theorem}

\begin{theorem}[Effects of new-frame code]
\lean{Light.Compiler.effects_newFrame}
\label{Light.Compiler.effects_newFrame}
\leanok
\uses{Light.Compiler.Framed, Light.Compiler.Rel, Light.Compiler.Rel.beyond, Light.Compiler.Rel.of_agreeOutside, Light.Compiler.Scratch, Light.Compiler.Sizes.inRange_stack, Light.Compiler.Stack, Light.Compiler.Writable, Light.Compiler.cADDR, Light.Compiler.cFP, Light.Compiler.cNFP, Light.Compiler.cStack, Light.Compiler.cStack_sub, Light.Compiler.cStack_succ_add, Light.Compiler.cT, Light.Compiler.effects_frameInit, Light.Compiler.frameInit, Light.Compiler.newFrame, Light.Compiler.scratch_subset_writable, Light.Compiler.stack_subset_writable, Light.Expr, Light.Expr.val, Light.frame, WordRAM.AgreeOutside, WordRAM.AgreeOutside.mono, WordRAM.AgreeOutside.read, WordRAM.AgreeOutside.refl, WordRAM.AgreeOutside.trans, WordRAM.AgreeOutside.update, WordRAM.effect_add, WordRAM.effect_store, WordRAM.effect_sub, WordRAM.effects, WordRAM.effects_append, WordRAM.effects_cons, WordRAM.effects_nil, WordRAM.wd}
\difficulty{5}
Let a word size and compiler size data specify a layout, limits, and stack capacity, and
let a state, stack position $q$, and memory satisfy the framed-state conditions. Suppose
$q+F+1+F\le Q$. For a list of argument expressions, assume each compiled argument value
is already in its corresponding temporary cell, and let the return position $r$ be at
most the dispatcher position. After executing the new-frame code, the memory and state
satisfy the relation for a new frame whose local variables are the argument values (with
missing arguments interpreted as zero), whose ordinary memory is unchanged, and whose
stack position is $q+F+1$. The stack cell at index $q+F$ contains $r$, and memory is
unchanged outside the writable locations for the new frame.
\end{theorem}
